\section{II lab: Hacking Team}

Hacking Team was an Italian technology company known for selling offensive intrusion and surveillance services to government and law enforcement entities worldwide. Founded in 2003, the firm became a highly controversial actor in the cybersecurity global landscape due to its commercial practices and specific clients



\subsection{Company Overview and Historical Evolution}
The company was established in Milan, Italy, by David Vincenzetti and Valeriano Bedeschi. 
\begin{itemize}
    \item \textbf{Foundation:} The company began operations, focusing its business model on providing offensive hacking tools and surveillance software to governments.
    \item \textbf{Early Success (2007):} Hacking Team received 8 million euros in funding from Italian venture capital firms. During this initial period, its software was purchased by the Milan police, allowing the company to establish itself as a commercial provider of hacking software
    \item \textbf{Global Reach:} Over its operational history, the company established high-level connections with major intelligence and defense actors globally, including Booz Allen Hamilton (specifically establishing relations with Mike McConnell, former NSA director and intelligence advisor) and United States intelligence agencies such as the NSA, CIA, and FBI.
    \item \textbf{Financial Growth:} The company grew significantly reporting revenues exceeding 40 million euros
\end{itemize}



\subsection{Core Technologies}
The technology developed by the company focused on \textbf{bypass-oriented} intrusion vectors designed to monitor communications directly from target devices. \\
From 2003 to 2015, the company achieved significant technical results in an environment completely lacking specific legislative frameworks.
\medskip

\noindent
The product relies on three  capabilities:
\begin{itemize}
    \item \textbf{Remote Control System (RCS):} Software designed for monitoring internet communications and decrypting target files directly at the endpoint
    \item \textbf{Mobile Surveillance:} Tools for tracking cellular phone metadata, call logs, active geographic location coordinates, and message contents
    \item \textbf{Remote Activation (Stealth Controls):} The capability to remotely activate microphones and webcams on target devices
\end{itemize}

\noindent
To ensure persistence and hide activities from security software, the platform utilized advanced surveillance techniques focused on operational stealth:
\begin{itemize}
    \item \textbf{Battery Optimization:} Techniques that allow to minimize processing overhead on monitored devices, avoiding suspicious battery loss
    \item \textbf{Stealth Operations:} Anti-detection methods designed to evade endpoint antivirus software and host-based intrusion detection systems (HIDS).
    \item \textbf{Data Extraction Tools:} Modules to extract data from various applications, such as email, database ecc...
\end{itemize}



\subsection{Legal Countermeasures}

As Hacking Team expanded its operations, its commercial decisions prompted intense international responses, regulatory interventions, and formal investigations.

\subsubsection*{Saudi Arabia Acquisition }
In 2013, the company became the subject of secret acquisition for the government of Saudi Arabia.\\
During the negotiations, the company was valued at approximately 2 billion dollars.\\ Saudi Arabian government extended a formal purchase offer of 140 million dollars. This event intensified external destination and potential misuse of the technology.

\subsubsection*{USA investigation about Sudan}
The most significant international regulatory involved the sale of software to Sudan:
\begin{itemize}
    \item \textbf{June 2014:} A United Nations commission initiated an official inquiry regarding software sales executed by Hacking Team to the Sudan government, given the active UN arms embargo over the region.
    \item \textbf{January 2015:} Hacking Team formally denied ongoing software sales to Sudan.
    \item \textbf{March 2015:} The UN commission insisted that the software could be classified as military equipment due to its offensive features and usage parameters.
\end{itemize}


\subsubsection*{Italian Ban}
In the autumn of 2014, the Italian government block Hacking Team's export capabilities. 
\begin{itemize}
    \item From a defense standpoint, while offensive malware is recognized as a necessary asset for combating cybercrime, exporting such powerful must be regulated to prevent overriding international human rights.
    \item From law standpoint, blocking commercial can conflict with constitutional rights regarding corporate competition and economic enterprise.
\end{itemize}


\subsection{Data Breach and Global Consequences}
Hacking Team in 2015 suffered a major data breach

\subsubsection*{The WikiLeaks Leak and Exposure of Confidential Data}
The breach exfiltratrated approximately 200 / 400 G of confidential data.\\
These data was published online on repositories such as WikiLeaks and a dedicated "Hacking Team Wiki" page.
\medskip

\noindent
The leak exposed the complete client list of the company.\\
So they revealed that Hacking Team had sold its software to highly restrictive regimes, including North Korea and South Korea.\\ Furthermore, documentation showed that the software had been sold to the DEA (the United States Drug Enforcement Administration) and simultaneously to a known narcotics trafficker in Mexico.
\medskip

\noindent
The leak also verified a strong connection between Hacking Team software deployments and Egypt, notably overlapping with the timeline of the Giulio Regeni case.

\subsubsection*{Impacts}
The real-world consequences of the data leak were immediate. \\
Within seven days of the leak's publication, South Korean spy and a North Korea spy whose activities were logged within the leaked email.\\
Moreover, more than 20 individuals died within a single month following the breach due to the exposure of data names, investigative logs, and target lists contained in the leaked files.


\subsection{Investigation and Judicial Realities}

\subsubsection*{The SoftHack S.r.l. Search and Seizure}
Prosecutor Alessandro Gobbis (Milan procetutor), initiated an official criminal investigation triggered by a highly suspicious financial payment from a Saudi Arabian company to a Turin cybersecurity firm named \textit{SoftHack S.r.l}. 
\begin{itemize}
    \item \textbf{The Investigative Trigger:} SoftHack S.r.l. had been founded by a group of former employees who had left Hacking Team.
    \item \textbf{The Corporate Theory:} David Vincenzetti suspected that these employees had stolen Hacking Team's  source code of the \textit{Galileo} spyware platform (the core engine of the RCS system).\\
    Vincenzetti's theory was that these ex-employees had created SoftHack to illicitly sell Hacking Team's stolen surveillance tools to unauthorized external entities and unregulated foreign states, such as Saudi Arabia.
    \item \textbf{Prosecutorial Action:} Based on the suspicious payment trail the prosecutor ordered a formal search and seizure operation against the Turin company to investigate the potential unauthorized sale of the spyware source code.
\end{itemize}
Despite years of intense global controversy and public scrutiny over its choice of clients, Hacking Team itself was never formally indicted or prosecuted as an entity for its surveillance sales. The national judicial system was only mobilized when the issue shifted into a dispute over corporate trade secrets and the alleged internal theft of malware code.

\subsubsection*{Arrest of David Vincenzetti}
Vincenzetti was arrested because he maintained a large collection of historical weapons at his private residence and, during a violent domestic dispute, attempted to attack his cousin with a sword
\medskip

\noindent
Following his arrest, defense counsel representing the accused SoftHack employees utilized Vincenzetti's personal behavior to outline a psychological profile for the prosecutor.\\
The defense won. They said the trade-secret probe of SoftHack S.r.l. was started by a very unstable man who sued former employees harshly just because they left his company


\subsection{Legal Implications}

The operational history of Hacking Team between 2003 and 2015 highlights a severe legislative vacuum regarding remote computer network exploitation tools.\\
\medskip

\noindent
From a constitutional law standpoint, treating a remote trojan injection identically to standard telecommunications interception represents a severe legal error:
\begin{itemize}
    \item \textbf{Traditional Telecommunications Interception:} Involves capturing a live conversation or data stream in transit over a specific, limited time window (traditionally capped at a maximum of 90 days under the Italian legal system).
    \item \textbf{Remote Trojan Executio:} Bypasses the transit layer to execute a thorough search of the local device. By accessing a user's persistent data archives (whatsapp chat).
\end{itemize}
